\section{Introduction}
\label{sec:introduction}

Video is increasingly the primary medium for consumer content,
education, and public communication. As audiences become more global,
the demand for localized versions of the same video has grown;
yet while dubbing and subtitle generation are now routine, the
\emph{scene text} that appears \emph{inside} the video---signs,
menus, banners, product labels, handwritten notes captured in the
frame---remains largely untouched by automated localization. Manual
alternatives such as rotoscoping, frame-by-frame repainting, or
3D tracking with replacement assets are labor-intensive and do not
scale. Automating the process of \emph{cross-language scene text
replacement}, in which source-language text is replaced in place with
a target-language equivalent, is therefore an attractive engineering
goal.

The task is notoriously difficult. A usable replacement must be
simultaneously (i)~\emph{translated} correctly into the target
language, (ii)~\emph{geometrically aligned} with the original text
quadrilateral in every frame, (iii)~\emph{temporally consistent} so
that the result does not flicker or drift across frames, and
(iv)~\emph{visually realistic}, inheriting the style, texture,
lighting, and motion blur of the surrounding scene. A failure on any
of these axes is immediately visible to a viewer.

Two complementary paradigms currently address this space. The first
is \emph{end-to-end generative video editing}: systems such as
Sora~\citep{openai2024sora} and Runway's Gen-4 Aleph~\citep{runway2025aleph}
rewrite video frames at unprecedented photorealistic quality, and
their central strength is naturalism---lighting, texture, and motion
are preserved in a way that hand-crafted compositors cannot currently
match. Their limitation, at the time of writing, is \emph{glyph-level
control}: when prompted for a specific translation, these models
frequently produce unreadable or incorrect target-language strings,
because reliably rendering glyphs conditioned on a target string
remains an open research problem. The second paradigm is
\emph{glyph-controlled compositing}: an explicit pipeline that
detects existing text, generates the target string in an editor that
accepts the string as input, and composites the result back into the
video. This paradigm trades some of the naturalism of generative
systems for predictable control over what text appears on screen.

Our work occupies the second paradigm. We present a prototype system
for cross-language scene text replacement that takes as input a
source-language video and a target language, and produces a video
with the same content but with all detected scene text replaced by
correct translations. The system is structured as a five-stage
pipeline integrating seven machine-learning models, of which two we
implemented and trained ourselves: an Alignment Refiner of our own
design that addresses residual tracking error in the compositing
path, and a Blur Prediction Network (BPN) following STRIVE's
formulation~\citep{subramanian2021strive} that handles per-frame
blur adaptation. STRIVE's code is unreleased, so our BPN is a
re-implementation rather than a design from scratch; we differ from
STRIVE in backbone choice, training data preparation, and one
inference-time robustness fix, detailed in
Section~\ref{sec:method:bpn}. The remaining five components
(PaddleOCR, CoTracker3, AnyText2, SRNet, and Hi-SAM) are pretrained
and used through stable interfaces. The
whole pipeline is packaged as a reproducible open reference
implementation.

We view this system as a \emph{strong baseline} in the
glyph-controlled paradigm rather than a competitor to end-to-end
generative video editing. The generative paradigm is currently
winning on style and is likely to dominate as glyph control matures
in those systems; until then, modular pipelines such as ours offer
the predictable correctness that end-to-end generative models cannot
yet guarantee. The two approaches are complementary, and we expect
that hybrid systems combining generative style with glyph-controlled
compositing are the most promising long-term direction.

The remainder of the report is structured as follows.
Section~\ref{sec:related} surveys related work in image scene text
editing, video text replacement, generative video editing, and the
supporting vision components we build on.
Section~\ref{sec:method} presents the full pipeline stage by stage
and gives detailed accounts of the two models we trained.
Section~\ref{sec:experiments} presents qualitative comparisons
against a generative baseline, per-module ablations, and a brief
discussion of current limitations. Section~\ref{sec:conclusion}
closes with an honest assessment of the system's contribution and
open directions for future work.

\figplaceholder{hero_webui}{\linewidth}%
  {Our prototype end-to-end system, shown via its web interface:
   the user uploads a video, selects source and target languages,
   and receives a localized output in which scene text has been
   replaced frame by frame. Internally the system integrates seven
   ML models across five stages.}%
  {fig:hero}
